\section{積分の不等式}\label{節:積分の不等式}

$f(x)\le g(x)$が区間全体で成り立てば,\,$\int_a^b f(x)\dd x\le\int_a^b g(x)\dd x$である.この基本的な性質を積や平方の形に使うと,2つの関数の積分の間に成り立つ不等式が得られる.

\begin{theorem}{積分に対するコーシー・シュワルツの不等式}
$a<b$とし,\,$f,g$は$[a,b]$で実数値を取る連続関数とする.このとき
\begin{equation}\label{式:積分のコーシーシュワルツ}
 \left(\int_a^b f(x)g(x)\dd x\right)^2
 \le
 \left(\int_a^b f(x)^2\dd x\right)
 \left(\int_a^b g(x)^2\dd x\right)
\end{equation}
が成り立つ.\,$g$が恒等的に$0$でなければ,等号が成り立つための必要十分条件は,ある定数$c$について区間全体で$f(x)=cg(x)$となることである.\,$g$が恒等的に$0$なら,どの$f$についても等号が成り立つ.
\end{theorem}

平方は負にならないので,任意の実数$t$について
\[
 0\le\int_a^b\{f(x)-tg(x)\}^2\dd x
\]
である.式を短くするため
\[
 A=\int_a^b f(x)^2\dd x,\qquad
 B=\int_a^b f(x)g(x)\dd x,\qquad
 C=\int_a^b g(x)^2\dd x
\]
と置くと,右辺を展開して
\[
 0\le Ct^2-2Bt+A
\]
となる.

$C>0$とする\footnote{$C=0$なら,下の等号の議論と同じく連続性から$g$は区間全体で$0$となり,\,$B=0$である.このとき不等式の両辺はともに$0$である.}.右辺は$t$の2次関数であり,すべての実数$t$で$0$以上である.ここで1つの知識を思い出そう.2次関数がすべての実数で$0$以上なら,その判別式は$0$以下である.右辺の判別式を$D$とすると
\[
 \frac{D}{4}=B^2-AC\le0
\]
となり,\,$B^2\le AC$を得る.これが式\eqref{式:積分のコーシーシュワルツ}である.

等号$B^2=AC$が成り立つのは判別式が$0$のときであり,そのとき2次関数は$t=\frac BC$で$0$になる.つまり
\[
 \int_a^b\left\{f(x)-\frac BCg(x)\right\}^2\dd x=0.
\]
被積分関数は連続かつ$0$以上である.もし1点で正なら,連続性によってその近くの小区間でも正となり,積分が正になってしまう.したがって区間全体で$f(x)=\frac BCg(x)$である.逆に$f=cg$なら,直接代入して等号を確かめられる.

この不等式を使うときには,評価したい積分の中身を$f(x)g(x)$という積に分ける.片方を$1$とすることも多い.例えば$g(x)=1$とすれば
\begin{equation}\label{式:平均と二乗平均}
 \left(\int_a^b f(x)\dd x\right)^2
 \le(b-a)\int_a^b f(x)^2\dd x
\end{equation}
となる.両辺を$(b-a)^2$で割ると,関数の平均値の平方は,平方の平均値を超えないことが分かる.等号は$f$が定数関数のときに限る.

\begin{bookpoint}
\textbf{POINT}\par
評価したい積分の中身を$f(x)g(x)$の積に分け,\,$f^2$と$g^2$の積分が求めたい形になるように$f,g$を選んで式\eqref{式:積分のコーシーシュワルツ}に当てはめる.
\end{bookpoint}

\begin{example}{逆数と不等式}
$h$は$[a,b]$で連続で,すべての$x$について$h(x)>0$を満たすとする.
\[
 \left(\int_a^b h(x)\dd x\right)
 \left(\int_a^b\frac1{h(x)}\dd x\right)
 \ge(b-a)^2
\]
を示し,等号条件を求めよ.
\end{example}

\paragraphbf{方針}
2つの積分を右辺の$\int f^2$と$\int g^2$にするため,\,$f=\sqrt h$,\,$g=\frac1{\sqrt h}$と置く.積$fg$は$1$になる.

\paragraphbf{解答}
不等式は成り立ち,等号は$h$が正の定数関数のとき,またそのときに限り成り立つ.

\paragraphbf{解法}
$f=\sqrt h$,\,$g=\frac1{\sqrt h}$と置く.\,$h$は閉区間で連続かつ正なので,この2つの関数も連続である.\,$f^2=h$,\,$g^2=\frac1h$,\,$fg=1$だから,式\eqref{式:積分のコーシーシュワルツ}より
\[
 (b-a)^2
 =\left(\int_a^b1\dd x\right)^2
 \le\left(\int_a^b h(x)\dd x\right)
       \left(\int_a^b\frac1{h(x)}\dd x\right)
\]
となる.\,$g$は恒等的に$0$ではないので,等号条件は$f=cg$である.これは$\sqrt h=\frac c{\sqrt h}$,すなわち$h=c$を意味する.

\paragraphbf{対数の評価}
例えば$h(x)=1+x$,\,$0\le x\le1$とすると,
\[
 \int_0^1(1+x)\dd x=\frac32,\qquad
 \int_0^1\frac1{1+x}\dd x=\log2
\]
である.\,$h$は定数ではないから等号は成り立たず,
\[
 \frac32\log2>1
\]
となる.つまり$\log2>\frac23=0.666\ldots$である.実際$\log2=0.693\ldots$であり,対数の値を小数で計算せずに,厳密な下からの評価が得られたのだ.
